\chapter[Compilación de OpenVDB]{Instrucciones de compilación de OpenVDB para Python}
\label{anx:compilacion_openvdb}
Este anexo presenta las instrucciones de compilación de la biblioteca OpenVDB para Python en 
Windows 11.

\section{Requisitos}
\begin{itemize}
    \item CMake
    \item Visual Studio 2022 con los siguientes componentes:
    \begin{itemize}
        \item Carga de trabajo <<Desarrollo para el escritorio con C++>>
        \item Herramientas de compilación de C++ de MSVC v143-VS 2022 para x64/x86 
        (más recientes)
        \item SDK de Windows 11
        \item Herramientas de CMake en C++ para Windows
        \item ATL de C++ para las herramientas de compilación v143 más recientes (x86 y x64)
    \end{itemize}
    \item Python
    \item Git
\end{itemize}
Los componentes de Visual Studio se pueden instalar mediante Visual Studio Installer, 
seleccionando la opción \textit{Modificar} en la versión del programa instalada. En la 
ventana emergente se debe marcar la carga de trabajo 
\textit{Desarrollo para el escritorio con C++}, y en el panel lateral derecho seleccionar 
los componentes opcionales requeridos. Adicionalmente, estos pueden localizarse de forma 
individual en la pestaña \textit{Componentes individuales}.

\begin{figure}[!htbp]
    \centering
    \includegraphics[width=0.85\textwidth]{images/anexos/compilacion-openvdb/vs-componentes.png}
    \caption[Instalación de componentes de Visual Studio]{Instalación de componentes de Visual 
    Studio.}
    \label{fig:vs_components}
\end{figure}

\section{Instrucciones}
La compilación de OpenVDB se realiza a través de la consola de comandos de Windows 
(CMD o PowerShell). En la terminal, ejecute los siguientes comandos para clonar y preparar el 
gestor de paquetes \textit{vcpkg} en el directorio raíz \texttt{C:\textbackslash}:

\begin{verbatim}
    cd C:\
    git clone https://github.com/microsoft/vcpkg.git
    cd vcpkg
    .\bootstrap-vcpkg.bat
\end{verbatim}

\textit{vcpkg} es un gestor de paquetes para C/C++ que compila e instala las dependencias 
necesarias para la biblioteca. Instale dichas dependencias mediante el siguiente comando:

\begin{verbatim}
    .\vcpkg install tbb blosc boost-iostreams boost-system boost-any 
    boost-algorithm openexr imath zlib nanobind boost-interprocess 
    boost-headers:x64-windows
\end{verbatim}
El proceso de descarga y compilación de dependencias puede demorar entre 15 y 20 minutos.
\medskip

Posteriormente, clone el código fuente de OpenVDB en el directorio raíz y cree el subdirectorio 
de compilación (\texttt{build}):

\begin{verbatim}
    cd C:\
    git clone https://github.com/AcademySoftwareFoundation/openvdb.git
    cd openvdb
    mkdir build
    cd build
\end{verbatim}

Configure el proyecto de compilación utilizando CMake con la siguiente instrucción (una sola linea):

\begin{verbatim}
    cmake .. -G "Visual Studio 17 2022" -A x64 ^
     -DCMAKE_PREFIX_PATH="C:/vcpkg/installed/x64-windows" ^
     -DOPENVDB_BUILD_CORE=ON -DOPENVDB_BUILD_PYTHON_MODULE=ON ^
     -DPython_EXECUTABLE="C:/path/a/instalacion/de/python/python.exe" ^
     -DUSE_BLOSC=ON -DUSE_TBB=ON
\end{verbatim}
Reemplace el valor asignado a \texttt{-DPython\_EXECUTABLE} por la ruta absoluta al ejecutable 
de Python (\texttt{python.exe}) correspondiente a la versión específica para la cual se 
compilará OpenVDB. Se recomienda utilizar la instalación nativa del sistema 
(y no el intérprete de un entorno virtual) para prevenir fallos en la resolución de bibliotecas 
dinámicas. La ruta estándar suele ubicarse en:

\begin{verbatim}
    C:\Users\<usuario>\AppData\Local\Programs\Python\Python3XX\python.exe
\end{verbatim}

Inicie el proceso de compilación ejecutando:

\begin{verbatim}
    cmake --build . --config Release --target openvdb_python
\end{verbatim}

Una vez finalizada la compilación, el módulo compilado se localizará en:

\begin{verbatim}
    C:\openvdb\build\openvdb\openvdb\python\Release\
\end{verbatim}
con el nombre de archivo \texttt{openvdb.cp3XX-win\_amd64.pyd} 
(donde \texttt{XX} corresponde a la versión menor de Python instalada).

\section{Instalación}
Para habilitar el módulo compilado en un entorno de Python, copie los siguientes archivos 
generados:

\begin{verbatim}
    Desde C:\openvdb\build\openvdb\openvdb\python\Release:
    - openvdb.cp3XX-win_amd64.pyd

    Desde C:\openvdb\build\openvdb\openvdb\Release:
    - openvdb.dll

    Desde C:\vcpkg\installed\x64-windows\bin:
    - tbb12.dll
    - zlib1.dll o z.dll
    - zstd.dll
    - blosc.dll
    - lz4.dll
\end{verbatim}
al directorio \texttt{Lib/site-packages} del entorno virtual (\texttt{.venv}) de destino, 
o al directorio análogo de la instalación global de Python.
\medskip

Para comprobar que el módulo fue enlazado correctamente, abra una sesión del intérprete 
interactivo de Python e importe el paquete:

\begin{verbatim}
    import openvdb
\end{verbatim}
Si la instrucción finaliza sin arrojar excepciones de tipo \texttt{ImportError} o 
fallos de enlace de librerías dinámicas DLL, la instalación habrá concluido de forma exitosa.